% tools/textfix/fixtures/sample.tex
% 测试样例：数字之间的全角冒号应被改为半角；中文冒号与 LaTeX 命令保持不变。
\section{时刻与时间间隔}

如图为某列火车的订单。订单详情中的时间 19：26、12：21 和“历时 16 时 55 分”含义是否相同？
\hxansblock{不相同。19：26、12：21 表示时刻；“历时 16 时 55 分”表示时间间隔。}

\begin{enumerate}
  \item 上课时间 8：00 为时刻。\hxparen{√}
  \item 8：30 上课，其中“8：30”指的是时间间隔。\hxparen{×}
  \item 某次列车运行 8：00：00 后到站。\hxparen{√}
  \item 两数之比为 1：2。\hxparen{×}
  \item 注意：这里的冒号应保持全角。
  \item 时间是：8 点，而 8：00 是时刻。
  \item 数学模式中的 $t=8：00$ 也应修正。
  \item 位移—时间图像记作 $x$-$t$ 图像。
  \item 位移大小为 $\dfrac{3}{2}\pi R$。
  \item \tkanswer{√}\hxparen{×} 等命令不受影响。
\end{enumerate}
